\section{The Magisterial Exposition: Catechism and \emph{Veritatis splendor}}

The Code uses the taxonomy; two documents of the 1990s expound it. The Catechism of the Catholic Church (1992; Latin typical edition 1997) gives natural law its catechetical statement in seven paragraphs, and John Paul II's encyclical \emph{Veritatis splendor} (6 August 1993) gives it the fullest doctrinal treatment of the modern papal magisterium. Neither is legislation; both are teaching documents whose weight is that of the ordinary papal magisterium in their respective genres, and both, examined closely, teach the \emph{classical} doctrine---participation, not intuitionism---with its scholastic architecture intact.\endnote{CCC 1954--1960, quoted below from the Holy See's current English web presentation of the Catechism with the Latin archive page of the same article as the governing text, both fetched, hashed, and read 2026-07-25; \emph{Veritatis splendor} nn.~12 and 40--45, quoted from the Holy See's official English web delivery, fetched, hashed, and read the same day (the Latin text in AAS 85 [1993] 1133--1228 is authoritative and was not collated). Paragraph numbers follow the typical edition and the encyclical's own numbering.}

\subsection{CCC 1950--1953: the taxonomy in one paragraph}

The article on the moral law opens, before reaching natural law, with a compact restatement of this entire study. Paragraph 1951 defines: \latin{Lex regula est agendi ab auctoritate competenti propter bonum commune promulgata}---law is a rule of acting, promulgated by competent authority for the common good---the four clauses of q.~90 a.~4 in catechetical order; and grounds: \latin{Omnis lex in Lege aeterna suam primam et ultimam invenit veritatem}---every law finds its first and ultimate truth in the eternal law---with Leo XIII's \latin{Ordinatio rationis lex nominatur} in the note. Then paragraph 1952 draws the map entire: the expressions of the moral law are diverse and all coherent among themselves---\latin{Lex aeterna, fons omnium legum in Deo; lex naturalis; Lex revelata Legem veterem Legemque continens novam seu evangelicam; tandem leges civiles et ecclesiasticae}---the eternal law, source in God of all laws; the natural law; the revealed Law, comprising Old and New; and finally civil and ecclesiastical laws. Paragraph 1953 adds the Christological seal: the moral law finds its fullness and unity in Christ, who ``is the end of the law'' (Rom 10:4).\endnote{CCC 1950--1953, quoted from the Latin delivery of the article (which carries the full article, 1950--1986) with the article's own working glosses; the English page registered for this study carries only the natural-moral-law section (1954--1960), so the English of 1950--1953 is not quoted. Paragraph 1951's note cites Leo XIII, \emph{Libertas praestantissimum}; 1952's list is the exact taxonomy of section 6, in the same order, with the eternal law placed as source rather than species---the Catechism teaching in one sentence what this article argued from q.~93.}

\subsection{CCC 1954--1960: the catechetical seven}

The Catechism's exposition of natural law is built almost entirely out of citations, and the choice of citations is the teaching. \textbf{1954}: man participates in the Creator's wisdom and goodness; the natural law expresses ``the original moral sense which enables man to discern by reason the good and the evil''---and the paragraph's centerpiece is Leo XIII, \emph{Libertas praestantissimum}: the natural law ``is written and engraved in the soul of each and every man, because it is human reason ordaining him to do good and forbidding him to sin''; but this prescription of human reason ``would not have the force of law if it were not the voice and interpreter of a higher reason to which our spirit and our freedom must be submitted.''\endnote{CCC 1954, with Leo XIII, \emph{Libertas praestantissimum} (1888) at its center (Latin: \latin{Ista vero humanae rationis praescriptio vim habere legis non potest, nisi quia altioris est vox atque interpres rationis}). The Leonine sentence is the anti-autonomy hinge: reason promulgates but does not originate. \emph{Veritatis splendor} n.~44 quotes the same passage at greater length (below).}

\textbf{1955} names the law with a deliberately doubled adjective---``the `divine and natural' law'' (\latin{Lex divina et naturalis})---shows man the way; its principal precepts ``are expressed in the Decalogue''; and it is called natural ``not in reference to the nature of irrational beings, but because reason which decrees it properly belongs to human nature,'' with Augustine (the eternal rules written ``in the book of that light we call the truth,'' impressed on the heart ``like a seal on a ring that passes onto wax, without leaving the ring'') and Aquinas (natural law ``nothing other than the light of understanding placed in us by God'') as the two witnesses.\endnote{CCC 1955, quoting Augustine, \emph{De Trinitate} XIV, and Aquinas, \emph{In duo praecepta caritatis}---the same Thomas text VS 12 and 40 quote. The doubled adjective and the Decalogue clause do the compressing work Gratian once did, but now with the distinction available behind it: the Decalogue \emph{expresses} natural law's principal precepts in divinely revealed positive form---two registers, one content.}

\textbf{1956--1958} give the three classical properties. Universality: present in the heart of each man and established by reason, binding all men, expressing the dignity of the person and determining ``the basis for his fundamental rights and duties''---with Cicero's \emph{De re publica} as witness (``true law: right reason \ldots\ diffused among all men \ldots\ immutable and eternal''). Variability of application: ``application of the natural law varies greatly''; it ``can demand reflection that takes account of various conditions of life according to places, times, and circumstances''---yet in the diversity of cultures it ``remains as a rule that binds men among themselves.'' Immutability: ``immutable and permanent throughout the variations of history''; even ``when it is rejected in its very principles, it cannot be destroyed or removed from the heart of man,'' with Augustine's \emph{Confessions} on the law ``that iniquity itself does not efface.''\endnote{CCC 1956--1958. The three paragraphs are qq.~94 aa.~4--6 in catechetical dress, including their realism: universality of principle, variability of application, indelibility of the core. The choice of a pagan witness (Cicero) for universality is itself doctrine---the law is knowable outside revelation.}

\textbf{1959--1960} give the two juridical corollaries this article has been tracing all along. The natural law provides ``the necessary basis for the civil law with which it is connected, whether by a reflection that draws conclusions from its principles, or by additions of a positive and juridical nature''---q.~95 a.~2's conclusions and determinations, cited nearly verbatim in a universal catechism. And the limit of unaided nature: the precepts of natural law ``are not perceived by everyone clearly and immediately''; sinful man needs grace and revelation so that moral and religious truths may be known ``by everyone with facility, with firm certainty and with no admixture of error''---Pius XII's formula from \emph{Humani generis}, itself descended from q.~91 a.~4's reasons for divine positive law.\endnote{CCC 1959--1960. The 1960 citation chain (Pius XII quoting the First Vatican Council's epistemology of revelation) closes the loop this article opened at c.~748: the same doctrine---natural knowability in principle, revelation's necessity in the human condition---stands behind the canon's pairing of a divine-law obligation to seek truth with an absolute prohibition of coerced faith.}

\subsection{\emph{Veritatis splendor} 40--45: participated theonomy}

The encyclical's second chapter confronts the late-twentieth-century proposals directly, and its treatment of law (nn.~35--53) turns on a coined term. The occasion is the claim of ``autonomous morality'': that human reason itself creates values and norms. The encyclical's answer neither concedes autonomy nor retreats to heteronomy---law imposed as an alien will---but names a third thing.

\textbf{n.~40} states both sides of the truth: the moral life calls for the creativity and originality of the person, yet ``reason draws its own truth and authority from the eternal law, which is none other than divine wisdom itself.'' There is a ``rightful autonomy'' of the practical reason: ``man possesses in himself his own law, received from the Creator.'' \emph{Nevertheless}, ``the autonomy of reason cannot mean that reason itself creates values and moral norms.''\endnote{VS 40, whose note anchors ``eternal law'' in Augustine and in ST I-II q.~93 a.~1, and quotes Aquinas's \emph{In duo praecepta caritatis} once more for natural law.} \textbf{n.~41} coins the term: a heteronomy that denied man's self-determination would contradict revelation; what revelation actually teaches is ``theonomy, or \emph{participated theonomy}, since man's free obedience to God's law effectively implies that human reason and human will participate in God's wisdom and providence.'' Law is ``an expression of divine wisdom: by submitting to the law, freedom submits to the truth of creation.''\endnote{VS 41. ``Participated theonomy'' is the encyclical's one-phrase summary of I-II q.~91 a.~2; the term is new, the doctrine is the treatise's.}

\textbf{nn.~42--43} re-derive the classical structure with the classical texts. Man's freedom, ``patterned on God's freedom,'' is not negated by obedience to divine law but only thereby ``abides in the truth''; the discernment of good and evil happens ``thanks to the light of natural reason, the reflection in man of the splendour of God's countenance''---and n.~42 quotes, at length and by name, Aquinas's Psalm-4 gloss from q.~91 a.~2, adding that the law is called natural ``because the reason which promulgates it is proper to human nature.'' n.~43 then quotes the Council---``the supreme rule of life is the divine law itself, the eternal, objective and universal law by which God out of his wisdom and love arranges, directs and governs the whole world and the paths of the human community''---identifies the doctrine as ``the classic teaching on God's eternal law,'' cites Augustine's definition and Thomas's, and concludes: ``the natural law enters here as the human expression of God's eternal law,'' quoting q.~91 a.~2's participation formula entire.\endnote{VS 42--43. The conciliar sentence is \emph{Dignitatis humanae} 3, which this article verified at the official English delivery of the declaration (fetched 2026-07-25, byte-identical to the previously registered dated state): ``the highest norm of human life is the divine law---eternal, objective and universal---whereby God orders, directs and governs the entire universe and all the ways of the human community by a plan conceived in wisdom and love.'' DH 3's own note cites ST I-II q.~91 a.~1 and q.~93: the Council's natural-law teaching is expressly footnoted to the treatise this article has been reading.}

\textbf{n.~44} makes the reception explicit as history: ``the Church has often made reference to the Thomistic doctrine of natural law, including it in her own teaching on morality,'' and reproduces the whole Leonine passage the Catechism excerpts, through its conclusion: ``the natural law \emph{is itself the eternal law}, implanted in beings endowed with reason \ldots\ it is none other than the eternal reason of the Creator and Ruler of the universe.'' \textbf{n.~45} completes the taxonomy from the top: the Church receives the New Law, ``the `fulfilment' of God's law in Jesus Christ,'' interior, ``the law of the Spirit of life in Christ Jesus''; and---in the encyclical's own summary of everything this article has mapped---``even if moral-theological reflection usually distinguishes between the positive or revealed law of God and the natural law, and, within the economy of salvation, between the `old' and the `new' law, it must not be forgotten that these and other useful distinctions always refer to that law whose author is the one and the same God and which is always meant for man.''\endnote{VS 44--45. n.~12, earlier in the encyclical, had already laid the same foundation exegetically: God answers the question about the good ``by creating man and ordering him with wisdom and love to his final end, through the law which is inscribed in his heart (cf.\ \emph{Rom} 2:15), the `natural law','' then in the Decalogue, then in the New Covenant's law written on hearts.}

The magisterial exposition, then, adds no fifth category and subtracts none. What it adds is a defensive precision the thirteenth century did not need: against autonomy, that reason promulgates but does not author; against heteronomy, that the author's law is not alien to the creature it constitutes; and against fragmentation, that the taxonomy's distinctions---natural and positive, old and new---are distinctions \emph{within one legislator's one wisdom}. That last sentence of n.~45 is, in effect, the eternal law of q.~93 restated as a warning to people who would play the categories against each other.
